% Nội dung mục 4.4 -- tạo từ KDDTrain+ và kiểm tra bằng Python.
% Không đặt \documentclass hay \begin{document} trong tệp thành phần.
\section{Kết quả khoảng tin cậy 95\% cho năm nhóm}
\label{sec:ketqua-wilson}

Theo quy trình tại Mục~\ref{sec:phuongphap-wilson}, Bảng~\ref{tab:anova:wilson} và
Hình~\ref{fig:anova:wilson} trình bày ước lượng Wilson từ toàn bộ KDDTrain+
($n=125.973$). Các giá trị mô tả cơ cấu của mẫu train, không phải tỷ lệ tấn công của
Internet hiện tại.

\begin{table}[H]
\centering
\caption{Ước lượng tỷ lệ và khoảng tin cậy Wilson 95\% trên KDDTrain+}
\label{tab:anova:wilson}
\small
\setlength{\tabcolsep}{6pt}
\begin{tabular}{lrrr}
\toprule
Nhóm & Số quan sát & Tỷ lệ (\%) & Khoảng tin cậy 95\% (\%) \\
\midrule
Normal & 67.343 & 53,458 & [53,183; 53,734] \\
DoS    & 45.927 & 36,458 & [36,192; 36,724] \\
Probe  & 11.656 &  9,253 & [ 9,094;  9,414] \\
R2L    &    995 &  0,790 & [ 0,742;  0,840] \\
U2R    &     52 &  0,041 & [ 0,031;  0,054] \\
\midrule
Tổng   & 125.973 & 100,000 & -- \\
\bottomrule
\end{tabular}
\end{table}

\begin{figure}[H]
\centering
\includegraphics[width=0.88\linewidth]{../ANOVA/figures/wilson_95ci_five_groups.pdf}
\caption{Ước lượng điểm và khoảng tin cậy Wilson 95\% cho năm nhóm. Trục tỷ lệ dùng
thang logarit để khoảng của R2L và U2R vẫn quan sát được; kích thước mẫu được ghi cạnh
mỗi điểm.}
\label{fig:anova:wilson}
\end{figure}

Normal chiếm tỷ trọng lớn nhất, 53,458\% (KTC 95\%: 53,183--53,734\%), tiếp đến là
DoS với 36,458\% (36,192--36,724\%) và Probe với 9,253\% (9,094--9,414\%). Hai nhóm
hiếm R2L và U2R chỉ chiếm lần lượt 0,790\% và 0,041\%. Khoảng của U2R rộng hơn theo tỷ
lệ tương đối vì chỉ có 52 quan sát. Sự mất cân bằng này báo trước độ bất định lớn hơn
cho các ước lượng và hậu kiểm liên quan đến U2R; các khoảng riêng lẻ không được dùng như
một kiểm định so sánh tỷ lệ.
\FloatBarrier

